\subsection{ORDER RELATION BETWEEN CARDINALS}

The relation ``X is equipotent to a subset of Y'' is equivalent to ``there exists an injection of X into Y''; it is also equivalent to the relation ``Card (X) is equipotent to a subset of Card (Y)'' (Chapter II, § 3, no. 8, Theorem 1).

THEOREM 1. The relation \(R\{\mathfrak{x}, \mathfrak{y}\}\) : `` \(\mathfrak{x}\) and \(\mathfrak{y}\) are cardinals and \(\mathfrak{x}\) is equipotent to a subset of \(\mathfrak{y}\) '' is a well-ordering relation (§2, no. 1).

Since \(\mathbf{R}\{\mathfrak{x},\mathfrak{x}\}\) is true for every cardinal \(\mathfrak{x}\) , what must be proved is that, for every set \(\mathbf{E}\) of cardinals the relation `` \(\mathfrak{x} \in \mathbf{E}\) and \(\mathfrak{y} \in \mathbf{E}\) and \(\mathbf{R}\{\mathfrak{x},\mathfrak{y}\}\) '' is a well-ordering relation on \(\mathbf{E}\) . Consider the set \(\mathbf{A} = \bigcup_{\mathfrak{r}\in \mathbf{E}}\mathfrak{x}\) .

Every cardinal \(\mathfrak{r} \in \mathbf{E}\) is then a subset of A. By §2, Theorem 1 there exists a well-ordering relation on A, which we shall denote by \(\mathfrak{r} \leqslant y\) , and every subset of A is equipotent to a segment of A (§2, no. 5, Theorem 3, Corollary \(x\) ). For every cardinal \(\mathfrak{r} \in \mathbf{E}\) consider the set of segments of A which are equipotent to \(\mathfrak{r}\) ; this set of segments is not empty and therefore has a least element (§2, no. 1, Proposition 2); let \(\varphi(\mathfrak{r})\) denote this least element. The relation

`` \(x \in E\) and \(y \in E\) and \(x\) is equipotent to a subset of \(y\) ''

is then equivalent to

\[
" \mathfrak {x} \in \mathrm{E} \text {   and   } \mathfrak {y} \in \mathrm{E} \text {   and   } \varphi (\mathfrak {x}) \subset \varphi (\mathfrak {y})".
\]

For clearly the second relation implies the first. Conversely, if \(\mathfrak{r}\) is equipotent to a subset of \(\varphi(\mathfrak{y})\) , we cannot have \(\varphi(\mathfrak{y}) \subset \varphi(\mathfrak{x})\) and \(\varphi(\mathfrak{y}) \neq \varphi(\mathfrak{x})\) ; otherwise there would exist a segment of \(\varphi(\mathfrak{y})\) equipotent to \(\mathcal{X}\) ( \(\S\) 2, no. 5, Theorem 3, Corollary 3), contrary to the definition of \(\varphi(\mathfrak{x})\) . Since the set of segments of A is well-ordered by inclusion ( \(\S\) 2, no. 1, Proposition 2), the theorem follows.

We shall denote the relation \(\mathbf{R}\{\mathfrak{x},\mathfrak{y}\}\) by \(\mathfrak{x}\leqslant \mathfrak{y}\) . A set \(\mathbf{X}\) is equipotent to a subset of a set \(\mathbf{Y}\) if and only if \(\operatorname{Card}(\mathbf{X})\leqslant \operatorname{Card}(\mathbf{Y})\) .

Clearly we have \(0 \leqslant \mathfrak{r}\) for every cardinal \(\mathfrak{r}\) , and \(1 \leqslant \mathfrak{r}\) for every cardinal \(\mathfrak{r} \neq 0\) .

COROLLARY 1. Given any two sets, one of them is equipotent to a subset of the other.

COROLLARY 2. Two sets each of which is equipotent to a subset of the other are equipotent.

Remark. Given any set A, there exists a set whose elements are the cardinals \(\operatorname{Card}(X)\) for all the subsets X of A, namely, the set of objects of the form \(\operatorname{Card}(X)\) for \(X \in \mathfrak{P}(A)\) (Chapter II, §1, no. 6). For every cardinal a the relation ``r is a cardinal and \(r \leqslant a\) '' is therefore collectivizing in r (Chapter II, §1, no. 4), because it is equivalent to the relation `` \(\mathfrak{f}\) is of the form Card (X) for \(X \subset \mathfrak{a}\) ''; the set of all \(\mathfrak{f}\) satisfying this relation is called the set of cardinals \(\leqslant \mathfrak{a}\) .

PROPOSITION 2. For every family \((\mathfrak{a}_i)_{i\in I}\) of cardinals, there exists a unique cardinal \(\mathfrak{b}\) such that \(\mathfrak{a}_i\leqslant \mathfrak{b}\) for all \(i\in I\) and such that every cardinal \(\mathfrak{c}\) for which \(\mathfrak{a}_i\leqslant \mathfrak{c}\) for all \(i\in I\) is \(\geqslant \mathfrak{b}\) .

There exists a set E containing all the sets \(a_{i}\) (e.g., the sum of these sets (Chapter II, §4, no. 8)), whence \(a_{i} \leqslant a = Card(E)\) for all \(i \in I\) . The set F of cardinals \(\leqslant a\) is well-ordered and contains all the \(a_{i}\) , and therefore the family \((a_{i})_{i \in I}\) has a least upper bound b in F. Let c be a cardinal \(\geqslant a_{i}\) for all \(i \in I\) ; if \(c < b \leqslant a\) , then \(c \in F\) , and the inequality \(a_{i} \leqslant c\) contradicts the definition of the least upper bound of the family \((a_{i})\) in the ordered set F; hence the result.

By abuse of language, the cardinal \(\mathfrak{b}\) is called the least upper bound of the family \((\mathfrak{a}_i)_{i\in I}\) and is denoted by \(\sup_{i\in I}a_i\) .

PROPOSITION 3. Let X and Y be sets. If there exists a surjection f of X onto Y, then Card (Y) ≤ Card (X).

For there exists a section \(s\) associated with \(f\) (Chapter II, §3, no. 8, Proposition 8) and \(s\) is an injection of Y into X.

